% answers_short.tex —— 2002 年·全国新课程卷·理科 · 简短答案速查
% 由 tools/make_short_answers.py 自 answers.tex 的 \gksolhead/\gkanswer 提取，
% 并与 2002/2002_数学真题_全国新课程理_答案要点.md 逐题比对后生成，勿手改（改源 MD 后重跑）。

\begin{gknotice}
\textbf{使用说明}：本节只给最终答案便于快速核对，完整推导见解析版。答案为编写者独立推导，\textbf{非官方标准答案}。
\end{gknotice}

\gkpart{选择题}
\begin{center}\begin{tabular}{ccccccccccccc}
\toprule
题号 & 1 & 2 & 3 & 4 & 5 & 6 & 7 & 8 & 9 & 10 & 11 & 12 \\
\midrule
答案 & D & C & B & D & C & B & B & A & D & D & B & C \\
\bottomrule
\end{tabular}\end{center}

\gkpart{填空题}
\begin{itemize}[leftmargin=2.6em]
  \item[\textbf{13.}] $(0,0),(1,1)$。
  \item[\textbf{14.}] $1$。
  \item[\textbf{15.}] $\dfrac{3\pi}{\ln2}$。
  \item[\textbf{16.}] $\dfrac72$。
\end{itemize}

\gkpart{解答题}
\begin{itemize}[leftmargin=2.6em]
  \item[\textbf{17.}] $-\dfrac{31\sqrt2}{50}$。
  \item[\textbf{18.}] (1) 甲：坐标见解析版；乙：$MN=\sqrt{a^2-\sqrt2a+1}$；(2) 甲：$\dfrac\pi6$；乙：$a=\dfrac{\sqrt2}{2}$；(3) 甲：无此小问；乙：$\alpha=\arccos\left(-\dfrac13\right)$
  \item[\textbf{19.}] (1) $\dfrac{21}{32}$；(2) 5 人
  \item[\textbf{20.}] (1) $l:y=-\dfrac{x}{x_1^2}+\dfrac2{x_1}-a$；(2) $x_2=2x_1-ax_1^2$；$0<x_2\leqslant\dfrac1a$，且当 $x_1<\dfrac1a$ 时 $x_1<x_2<\dfrac1a$
  \item[\textbf{21.}] (1) 圆 $x^2+y^2=3$ 的右半圆（$x>0$）；(2) $\tan\theta=|y_0|$
  \item[\textbf{22.}] (1) $a_3=2$；(2) $a_n=a_{n-2}+2$，$n=3,4,5,\cdots$；(3) $a_n=n+(-1)^n$；$S_n=\dfrac{n(n+1)}2+\dfrac{(-1)^n-1}{2}$
\end{itemize}
